% RESEARCH_DRAFT_SHA256: a703478d5b1c3e5c84cc86005ce9676485a7903012095f8885249423bb2c4843
% Layout and source-reference wrappers expanded for the web renderer.
\documentclass[11pt]{article}
\usepackage[margin=25mm]{geometry}
\usepackage{fontspec,unicode-math}
\setmainfont{Latin Modern Roman}
\setmathfont{Latin Modern Math}
\usepackage{amsmath,amssymb,mathtools,xurl,hyperref}
\hypersetup{colorlinks=true,urlcolor=blue,linkcolor=blue}
\setlength{\parindent}{0pt}
\setlength{\parskip}{5pt}
\setlength{\emergencystretch}{3em}
\newcommand{\sref}[2]{\hyperlink{source:#1:#2}{[S#2]}}
\newcommand{\cataloguescope}[1]{\paragraph{Source question.}#1}
\newcommand{\floor}[1]{\left\lfloor#1\right\rfloor}
\begin{document}

\section*{Polygonization simulation: admissibility and the limit of rotational descent}
\noindent catalogue entry; catalogue entry. Research attempt dated 11 October 2026.
\paragraph{Result in scope.}This manuscript contains a source-specific attempt and complete proofs of the stated auxiliary calculations. It makes no new solution or novelty claim for the entire target.
\section{Original problem and preserved definition}
The following definition and source list are reproduced from the repository. Any earlier inspection dates and reported literature claims in this preserved text describe that source record, not additional verification by this attempt.
% BEGIN PRESERVED DEFINITION
\subsection{Definitions and mathematical statement}

\paragraph{Shared notation.}Standard number systems are $\mathbb N,\mathbb Z,\mathbb Q,\mathbb R,\mathbb C$. Specialized notation and index ranges follow the cited source and the definitions below.


Dislocations are line defects in a crystal. Polygonization describes their organization into boundaries separating regions of slightly different lattice orientation. A numerical simulation approximates the stated dislocation-evolution equations.
\paragraph{Record scope and qualifications.}
The catalogue title is retained as metadata. Named certificates, contracts, bounds or frameworks in that title are not asserted as results of the original source extract.
The governing model and correctness or error criterion are not specified in the source question.

Source-record limitation: Simulation correctness for polygonization is requested without a mathematical model or error criterion.
\paragraph{Statement recorded in the supplied source.}
Can numerical methods correctly simulate polygonization in bent single-crystal beams, where dislocations align into polygonization patterns? \href{https://doi.org/10.4171/owr/2015/53}{[S2]}
\paragraph{Source status.}
The supplied catalogue labels this source record partially solved; that metadata was checked against the supplied snapshot on 2026-10-11. The label is preserved for provenance and is not an independent verification of present open status or mathematical correctness. The catalogue literature assessment is dated 2026-08-21.
\subsection{Short English statement}
Simulate the formation of polygonization patterns in bent crystals.
\subsection{Sources}
\begin{itemize}
\item[S1] Ulam AI and the UnsolvedMath Contributors, UnsolvedMath: A Curated Collection of Open Mathematics Problems, catalogue entry. Catalogue curation and metadata: CC BY 4.0. The source extract is retained; definitions, English summaries and qualifications are editorial additions.
\url{https://huggingface.co/datasets/ulamai/UnsolvedMath}.
\item[S2] Mini-Workshop: Scales in Plasticity (2016), Oberwolfach Reports, DOI 10.4171/owr/2015/53, pp. 3007--3048
\url{https://doi.org/10.4171/owr/2015/53}.
\item[S3] Polygonization in bent crystals
\url{https://doi.org/10.1016/j.ijplas.2013.02.010}.
\item[S4] Polygonization theory
\url{https://doi.org/10.1016/j.ijengsci.2012.03.005}.
\item[S5] Numerical polygonization simulations
\url{https://doi.org/10.1080/01418619708205157}.
\item Source-specific Oberwolfach Reports licences, including CC BY-SA 4.0 where applicable, are retained; see THIRD\_PARTY\_NOTICES.md.
\end{itemize}
% END PRESERVED DEFINITION
\subsection{Verbatim supplied catalogue question}
Can numerical methods correctly simulate polygonization in bent single-crystal beams, where dislocations align into polygonization patterns?

\section{Source model and limits of the question}
Blesgen's report, printed pages 3033--3034, specifies
$D\varphi=R_eU_eF_p$, with $R_e\in SO(3)$, and the additive slip ansatz
\[
F_p(\gamma)=I+\sum_{a=1}^I\gamma_a m_a n_a^T,
\qquad |m_a|=|n_a|=1,\quad m_a^Tn_a=0.
\]
Its incremental energy contains
$W_{\rm st}(R_e^TD\varphi F_p^{-1})$, angle-gradient energy, external work,
and regularized slip dissipation. It suggests improved treatment of the
micro-rotations because local minimization can select unwanted states.
The catalogue's request for a physically correct polygonization pattern
does not give a beam experiment, parameter values, convergence norm, or
acceptance criterion. We attempt the suggested rotational descent route and
check two prerequisites. This is not a numerical beam simulation.

\section{An explicit admissibility obstruction for additive slips}
One individual shear has determinant one: in a basis beginning with
$m,n$, the matrix $I+\gamma mn^T$ is triangular with diagonal all ones.
This does not extend to a sum of shears. Take
$m_1=e_1,n_1=e_2,m_2=e_2,n_2=e_1$. Then
\[
F_p=\begin{pmatrix}1&\gamma_1&0\\\gamma_2&1&0\\0&0&1\end{pmatrix},
\qquad \det F_p=1-\gamma_1\gamma_2.
\]
The two prescribed systems satisfy every displayed orthogonality and unit
condition, but $\gamma_1=\gamma_2=1$ gives a singular matrix. Consequently
the displayed inverse in the energy is undefined there. If orientation
preservation is required, the admissible region in this example is
$\gamma_1\gamma_2<1$. The report's ansatz by itself does not justify
unconstrained slip iterates. This diagnoses a prerequisite for a solver;
it does not show that an admissible continuum formulation is inconsistent.

\section{A feasible rotational update and its precise descent guarantee}
For a real skew-symmetric $K$, set $R(h)=R\exp(-hK)$ with $R\in SO(3)$.
The power series gives $\exp(-hK)^T=\exp(hK)$, and multiplication of the
commuting series gives $\exp(hK)\exp(-hK)=I$. Thus $R(h)$ is orthogonal.
Its determinant is a continuous function of $h$ valued in $\{-1,1\}$;
at zero it is one, so it is always one. This update preserves the
pointwise rotation constraint exactly.

Let $E$ be a twice continuously differentiable, finite-dimensional
discretized energy on a neighborhood of the trial curve. On the vector
space of skew matrices use the Frobenius inner product. There is a unique
skew $K$ representing the differential by
$DE(R)[RA]=\langle K,A\rangle$ for every skew $A$: this follows by
expanding that linear functional in an orthonormal basis. Hence
$e(h)=E(R\exp(-hK))$ satisfies $e'(0)=-\|K\|_F^2$.
If, on $[0,h]$, the additional bound
$|e''(s)|\le L\|K\|_F^2$ holds, integration twice gives
\[
e(h)\le e(0)-h\|K\|_F^2+\tfrac12Lh^2\|K\|_F^2.
\]
For $K\ne0$ this is strict descent whenever $0<h<2/L$ (all positive
$h$ within the interval if $L=0$). No uniform $L$ for the original
continuum energy is supplied by this calculation. Multiple nodal rotations
use the direct sum of these tangent spaces; prescribed boundary nodes
must be held fixed. Finite-element interpolation and slip admissibility
are separate requirements, not consequences of nodal feasibility.

\paragraph{Why feasible descent does not select the correct state.}
On the circle subgroup of rotations about $e_3$, consider the illustrative
smooth energy
\[
E(\theta)=1-\cos(2\theta)+\delta(1-\cos\theta),\qquad 0<\delta<4.
\]
Both summands are nonnegative, and $E(0)=0$ is a global minimum.
But $E'(\pi)=0$, $E''(\pi)=4-\delta>0$, and $E(\pi)=2\delta>0$.
Continuity of the second derivative makes $\pi$ a strict local minimum
on that subgroup, so a gradient method started there stays there.
The energy is a polynomial in $\cos\theta=R_{11}$ on this subgroup.
This is a complete toy obstruction to the inference from descent to
global selection; it is not a claim about the actual beam energy's wells.

\section{Model-assisted audit and remaining scope}
The matrix set $SO(3)$ is compact: it is closed by its defining polynomial
equalities and bounded because each column has unit length. The source's
discussion of failure of compactness must therefore not be paraphrased as
noncompactness of this matrix group. Euler-angle representatives can be
unbounded even for a fixed rotation, since adding $2\pi k$ leaves that
rotation unchanged; fields require additional functional compactness
arguments. None is established here. The audit checked the inverse in the
source energy, the distinction between one slip and multiple slips, the
sign of the tangent derivative, and the local nature of the toy minimum.
These elementary calculations are not claimed novel. No physical data,
numerical experiment, PDE convergence proof, or theorem prover was used.
The proposed descent mechanism still lacks a global or physically justified
selection principle and a validated polygonization error criterion.

\section{Process and contributions}
Mathematics and model-assisted adversarial self-audit: OpenAI gpt-6-astra, ultra reasoning, as dispatched for this research batch; exact serving revision not exposed. Codex (GPT-6; exact serving revision not exposed) materially assisted with orchestration and format review. No human mathematical authorship, mathematical edits, or independent human proof checking is asserted.
The coordinating Codex (GPT-6; exact serving revision not exposed) also provided a separate model-assisted written-proof review of this manuscript. This is not independent human verification.
Contribution prompt A was applied to the full definition and supplied source row, then prompt B was applied adversarially to the derived results. The item-specific audit above records the scope checks and rejected extrapolations. Local assembly and structural checks do not establish mathematical correctness. The source catalogue status is preserved in TOP\_PROBLEM and is separate from this attempt status.
\paragraph{Integration review.} Codex (GPT-6; exact serving revision not exposed) reviewed the written arguments and expanded layout and source-reference wrappers. Repository metadata and the production web renderer were checked using Python 3.13 and JavaScript with MathJax 3.2.2. This model-assisted review is not independent human proof verification; no PDF compilation or formal proof checking is asserted.

\section{Completion Estimate}
\noindent\textbf{COMPLETION: 0\%}
The proved material is a rederivation of known or auxiliary facts rather than a substantiated new advance on the original target. No validated beam simulation, convergence theorem, or physical pattern-selection criterion is obtained. The calculations expose admissibility and local-minimum gaps in a suggested rotational optimization route.
\section{References}
The complete original Sources subsection above is retained, including catalogue attribution and licence notices. The additional list records primary sources actually consulted for this attempt on 11 October 2026; uncited source links reproduced above were not all rechecked.
\begin{itemize}
\item Blesgen, Deformation patterning in finite-strain Cosserat plasticity, OWR 53/2015, pp. 3033--3034; full model and polygonization question read in the cached primary report. The web DOI endpoint was unavailable; no additional physical validation is claimed. \url{https://doi.org/10.4171/owr/2015/53}.
\end{itemize}
\section{Final status}
\textbf{Unresolved in this attempt.} No validated beam simulation, convergence theorem, or physical pattern-selection criterion is obtained. The calculations expose admissibility and local-minimum gaps in a suggested rotational optimization route.
\end{document}
